% Source: https://artofproblemsolving.com/community/c3h1441930_radical_help
% Stable identity: aops-topic-1441930-q01
\textbf{(a)} The radical sign acts as a grouping symbol, like parentheses: $\sqrt{16+9}$ means $(16+9)^{\dfrac{1}{2}}$, so the sum is computed first and $\sqrt{16+9}=\sqrt{25}=\boxed{5}$. The order of operations says what to do first when no grouping is shown; it does not allow splitting a grouped expression. The student's step $\sqrt{16+9}=\sqrt{16}+\sqrt9$ uses the false ``rule'' $(a+b)^{\dfrac{1}{2}}=a^{\dfrac{1}{2}}+b^{\dfrac{1}{2}}$: powers distribute over \emph{products}, $(ab)^{\dfrac{1}{2}}=a^{\dfrac{1}{2}}b^{\dfrac{1}{2}}$, but not over sums. Indeed $7^2=49\ne25$.

\textbf{(b)} Both sides are positive, so it suffices to compare squares:
\[
\left(\sqrt a+\sqrt b\right)^2=a+b+2\sqrt{ab}>a+b=\left(\sqrt{a+b}\right)^2,
\]
since $\sqrt{ab}>0$. Hence $\sqrt{a+b}<\sqrt a+\sqrt b$; in particular $5<7$. $\blacksquare$
